% Clase 22 --- La curva de resonancia con resistencia (§4.1).
% Tres valores de R sobre la misma curva: el pico vale eps0/R (finito) y se
% angosta al bajar R, pero las COLAS son las mismas --- ahí no importa R. Eso es
% lo que hace que un receptor con R chica sea selectivo.
\begin{center}
\begin{tikzpicture}[scale=1]

  \begin{axis}[
    fisfig, width=10.5cm, height=5.6cm,
    xlabel={$\omega$}, ylabel={$I_M$},
    xmin=0.05, xmax=2.6, ymin=0, ymax=5.4,
    xtick={1}, xticklabels={$\frac{1}{\sqrt{LC}}$}, ytick=\empty,
    legend style={at={(0.97,0.95)}, anchor=north east, draw=figgray!40,
                  fill=white, font=\footnotesize},
  ]
    \addplot[figred, smooth, samples=300, domain=0.05:2.6]
      {1/sqrt(0.2^2 + (x - 1/x)^2)};
    \addlegendentry{$R$ chica}
    \addplot[figamber, smooth, samples=300, domain=0.05:2.6]
      {1/sqrt(0.5^2 + (x - 1/x)^2)};
    \addlegendentry{$R$ media}
    \addplot[figblue, smooth, samples=300, domain=0.05:2.6]
      {1/sqrt(1.1^2 + (x - 1/x)^2)};
    \addlegendentry{$R$ grande}
    \draw[figaux] (axis cs:1,0) -- (axis cs:1,5.2);
    \node[figlblaux, figred, anchor=west] at (axis cs:1.06,4.9)
      {pico $= \dfrac{\varepsilon_0}{R}$};
    \node[figlblaux, anchor=south east] at (axis cs:2.55,1.15)
      {las colas no dependen de $R$};
  \end{axis}

  \node[figlbl, anchor=north, align=center] at (5.2,-1.75)
    {$I_M = \dfrac{\varepsilon_0}{\sqrt{R^2 +
      \left(\omega L - \dfrac{1}{\omega C}\right)^2}}$};

\end{tikzpicture}\\[0.5em]
{\small\itshape Con resistencia la amplitud ya no diverge: el máximo ocurre
cuando la parte imaginaria de $\hat Z$ se anula ---o sea otra vez en
$\omega = 1/\sqrt{LC}$, la misma frecuencia--- pero ahí queda
$\hat Z = R$ y la corriente vale $\varepsilon_0/R$, grande pero finita. La
familia de curvas muestra el efecto de $R$ en las dos direcciones: al bajarla el
pico sube \emph{y} se angosta, mientras que las colas ---donde domina la
reactancia--- son prácticamente las mismas para todos. Eso explica de un vistazo
para qué sirve la resonancia en tecnología: un receptor de radio o un celular es
un circuito con $R$ muy chica sintonizado a la frecuencia de la emisora, de modo
que amplifica muchísimo esa señal y casi nada las vecinas; un microondas emite
en la frecuencia que resuena con la molécula de agua; las lámparas para tratar
la ictericia concentran la emisión en la frecuencia útil, sin calentar de más al
bebé. Y también explica el riesgo: si $R$ es demasiado chica, en resonancia la
corriente puede crecer hasta quemar el circuito ---o hasta que aparezcan efectos
no lineales que acoten la amplitud---.}
\end{center}
